\documentclass[11pt,a4paper]{article}
\usepackage[T1]{fontenc}
\usepackage[utf8]{inputenc}
\usepackage[brazilian]{babel}
\usepackage{lmodern,microtype}
\usepackage[a4paper,margin=2.3cm]{geometry}
\usepackage{array,booktabs,longtable,tabularx}
\usepackage{enumitem}
\usepackage[table]{xcolor}
\usepackage{hyperref,xurl}
\hypersetup{colorlinks=true,linkcolor=blue!55!black,urlcolor=blue!55!black}
\setlength{\parindent}{0pt}
\setlength{\parskip}{5pt}
\setlist[enumerate]{leftmargin=1.9em,itemsep=2pt}
\begin{document}
\begin{center}
{\LARGE\bfseries Conteúdos das provas objetivas de Matemática --- IFMT\par}
\vspace{5pt}
{\large Editais 166/2023, 029/2024 e 003/2026\par}
\vspace{5pt}
{\small Consolidação de conteúdos programáticos oficiais e observações de cadernos aplicados\par}
{\small Conferência documental: 4 de novembro de 2026}
\end{center}
\tableofcontents\newpage
\section{Objeto e critérios de confiabilidade}
Este material tem o objetivo de servir como índice de estudo para concursos efetivos da carreira de Professor EBTT do IFMT, área Matemática. Distingue \textbf{(1)} conteúdos programáticos publicados nos anexos; \textbf{(2)} estrutura dos cadernos efetivamente aplicados; \textbf{(3)} observações de questões específicas que puderam ser verificadas.

Os programas de Matemática de 2023 e de 2024 foram localizados nos respectivos \textbf{Anexos III}, ambos com 20 questões previstas de conhecimentos específicos. \textbf{As sínteses abaixo são paráfrases analíticas, não transcrições integrais de editais.} Não confundir esta disciplina com prova de seleção de estudantes ou concurso de outro cargo.

\section{Estrutura das provas efetivamente aplicadas}
\begin{tabularx}{\textwidth}{p{.19\textwidth}p{.16\textwidth}X}
\toprule
\textbf{Edital}&\textbf{Total}&\textbf{Divisão verificada na capa do caderno de Matemática}\\\midrule
166/2023&60 questões&
1--10 Língua Portuguesa; 11--20 Fundamentos da Educação e Legislação;
21--30 Conhecimentos Gerais de Mato Grosso; 31--40 Informática Básica;
41--60 Matemática (conhecimentos específicos).\\
029/2024&50 questões&
1--10 Língua Portuguesa; 11--20 Fundamentos em Educação, Conhecimentos Transversais e Legislação;
21--30 Conhecimentos Gerais do Estado de Mato Grosso;
31--50 Matemática (conhecimentos específicos).\\
003/2026&Não transcrito&
A prova foi realizada em 30/08/2026, mas o caderno específico de Matemática e o Anexo III não foram acessados integralmente para validação neste levantamento; conferir página oficial antes de registrar distribuição ou tópicos.\\
\bottomrule
\end{tabularx}

\subsection{PDFs oficiais dos cadernos de 2023 e 2024}
\begin{itemize}
\item \textbf{166/2023 (Matemática, 27 páginas):} \url{https://processoseletivo.ifmt.edu.br/media/filer_public/ba/c8/bac8e6d4-4154-4915-8632-055ef2f67862/caderno-de-provas-matematica.pdf}.
\item \textbf{029/2024 (Matemática, 27 páginas):} \url{https://processoseletivo.ifmt.edu.br/media/filer_public/d1/56/d156aa75-a73c-4bf3-9a7f-8add0b68a36c/matematica_v08.pdf}.
\item \textbf{003/2026:} \url{https://seletivo.ifmt.edu.br/edital/visualizar/173/}. Caderno de Matemática ainda não identificado por URL verificável.
\end{itemize}

\section{Conhecimentos específicos de Matemática: núcleo comum de 2023 e 2024}
\textbf{Constatação documental:} os Anexos III de 166/2023 e de 029/2024 (consolidado após retificação) apresentam \textbf{17 itens de Matemática com a mesma estrutura e conteúdos}. Por isso, os reunimos sem duplicação; na revisão de um novo concurso, verificar o Anexo III específico da edição futura.

\begin{enumerate}
\item \textbf{Lógica matemática:} noções de proposições, conectivos e operações lógicas; equivalências e implicações; leitura de afirmações matemáticas.
\item \textbf{Conjuntos e funções:} definição e operações com conjuntos, tipos de funções, propriedades, representações, aplicações; equações, inequações e sistemas.
\item \textbf{Sequências e progressões:} sequências, progressões aritméticas e geométricas, termos, somas e aplicações.
\item \textbf{Matemática financeira e Estatística:} porcentagens e operações financeiras; organização e interpretação de tabelas e gráficos; medidas de centralidade e dispersão.
\item \textbf{Análise combinatória e Probabilidade:} princípio fundamental de contagem; arranjos, combinações, permutações; noções de probabilidade e aplicações.
\item \textbf{Matrizes, determinantes e sistemas lineares:} operações e propriedades, resolução de sistemas, classificações e aplicações.
\item \textbf{Geometria plana e espacial:} conceitos fundamentais; congruência, propriedades dos triângulos, Teorema de Tales; áreas, perímetros; posições relativas de retas e planos; áreas de superfícies e volumes de sólidos; unidades do Sistema Internacional.
\item \textbf{Geometria analítica plana e Trigonometria:} sistemas de coordenadas; distância entre pontos e de ponto a reta; cônicas; relações em triângulo retângulo e círculo trigonométrico; funções, equações e inequações trigonométricas.
\item \textbf{Limites e continuidade de funções reais:} propriedades de limites, continuidade, limites infinitos e limites no infinito para funções reais de uma variável.
\item \textbf{Derivadas e antiderivadas:} regras de derivação, derivadas sucessivas, derivação implícita e regra da cadeia; taxas relacionadas; crescimento e concavidade; máximos e mínimos locais; antidiferenciação e integral indefinida.
\item \textbf{Integração:} técnicas de integração, integral definida e aplicações.
\item \textbf{Álgebra Linear:} espaços vetoriais de dimensão finita; transformações lineares e Teorema do Núcleo e da Imagem.
\item \textbf{Equações Diferenciais Ordinárias:} equações de primeira e segunda ordem com coeficientes constantes.
\item \textbf{Didática da Matemática:} contrato didático e seus efeitos; obstáculos didáticos; avaliação; interpretação de situações de ensino e aprendizagem, concepções e erros de estudantes.
\item \textbf{Tecnologias e jogos no ensino:} uso intencional de tecnologia e jogos como procedimentos de ensino de Matemática.
\item \textbf{História, modelagem e resolução de problemas:} História da Matemática como recurso e objeto de estudo; modelagem matemática e metodologias de resolução de problemas.
\item \textbf{Base Nacional Comum Curricular:} princípios curriculares e habilidades da BNCC relevantes ao ensino de Matemática.
\end{enumerate}

\textbf{Fontes do núcleo comum:}
\begin{itemize}
\item Edital 166/2023, anexos consolidados, Anexo III, área Matemática: \url{https://seletivo.ifmt.edu.br/edital/arquivos/889/}.
\item Edital 029/2024, anexos consolidados após Retificador 01/2024, Anexo III, área Matemática: \url{https://processoseletivo.ifmt.edu.br/media/filer_public/2c/de/2cde91b5-8ae1-4067-8e0c-37e67430f39f/edital_2024_029__anexos_cp_ebtt_consolidado_01.pdf}.
\end{itemize}
\textbf{Cuidado:} o nível de detalhe desta enumeração é de organização para o estudo. A redação completa e vinculante é a do anexo oficial de cada certame. O fato de um tópico constar no programa não garante que uma questão tenha sido elaborada sobre ele.

\section{Demais disciplinas da prova: comparação e escopo}
\subsection{Língua Portuguesa: 10 questões em ambas as edições}
Os programas de 2023 e 2024 contemplam leitura e interpretação de gêneros discursivos, condições de produção, sequências textuais, linguagem verbal e não verbal, variação linguística, adequação de registro, funções da linguagem, morfologia, sintaxe, concordância, regência, colocação pronominal, semântica, coesão e coerência, figuras de linguagem, ortografia, pontuação e acentuação. \textbf{Estudo recomendado:} resolver as dez primeiras questões de cada caderno para distinguir gramática aplicada e interpretação de texto.

\subsection{Fundamentos em Educação, conhecimentos transversais e legislação: 10 questões}
Os anexos incluem função social da educação escolar, ensino e aprendizagem na Educação Básica e Profissional, currículo, planejamento e metodologia, tecnologias digitais na educação, avaliação institucional e da aprendizagem, organização do trabalho escolar e legislação aplicável.

No núcleo legislativo constam, entre outros, Constituição Federal, LDB (Lei 9.394/1996), criação dos Institutos Federais (Lei 11.892/2008), Regulamento Didático do IFMT, carreira EBTT (Lei 12.772/2012), regime jurídico dos servidores federais (Lei 8.112/1990), ética no serviço público, direitos humanos, LGPD, políticas de educação e inclusão e normas correlatas que aparecem nos programas históricos.

\textbf{Importante:} para legislações, conferir o texto e a vigência exigidos no edital correspondente; estes números de leis são referências históricas de cobrança, não substituem leitura dos anexos completos nem consolidação legislativa atual.

\subsection{Conhecimentos Gerais de Mato Grosso: 10 questões}
Temas como formação populacional, migrações e diversidade, indicadores sociais, rede urbana, geomorfologia, solos, bacias hidrográficas, clima, biomas, agricultura e pecuária, desenvolvimento regional, indústria, serviços, logística, turismo, fontes de energia, impactos ambientais, sustentabilidade e atualidades aparecem nos anexos. Organizar mapas, linhas do tempo e exercícios de interpretação do contexto do Estado.

\subsection{Informática: diferença importante entre 2023 e 2024}
A edição 166/2023 teve \textbf{10 questões próprias de Informática Básica}. Seu programa inclui ambiente Linux, Windows, diretórios e arquivos, softwares básicos de produtividade, noções de computação e tópicos associados. Já o caderno de Matemática 029/2024 não apresenta bloco autônomo de questões de informática: a faixa 31--50 é dedicada a conhecimentos específicos. Portanto, \textbf{não presumir que a matriz de 2023 se repetirá nas edições posteriores}.

\section{Exemplos de conteúdos identificados em questões efetivamente aplicadas}
A seguir estão apenas \textbf{exemplos verificados}, não uma classificação exaustiva das quarenta questões específicas das duas edições. Os enunciados originais devem ser consultados nos PDFs oficiais; não são reproduzidos aqui.
\begin{longtable}{p{.16\textwidth}p{.16\textwidth}p{.55\textwidth}}
\toprule
\textbf{Edital}&\textbf{Questão}&\textbf{Assunto matemático observado}\\\midrule\endhead
166/2023 & 41 & Interpretação de evolução de população e taxas crescentes/decrescentes; leitura de comportamento qualitativo de funções.\\
166/2023 & 43 & Regra da cadeia para derivada de uma função composta de polinômios.\\
166/2023 & 44 & Recuperação de grandeza a partir de taxa de variação, condição inicial e integral.\\
166/2023 & 53 & Triângulos equiláteros sucessivos, razões de progressões geométricas para lados, alturas e áreas.\\
029/2024 & 31 & Receita como produto de preço e quantidade, partindo de função de demanda afim.\\
029/2024 & 32 & Função definida por partes, extensão removível e continuidade/limites.\\
\bottomrule
\end{longtable}
\textbf{Atenção:} a classificação acima é didática e aproximada. Não atribuir ao edital um peso oficial por assunto a partir desta pequena amostra.

\section{Situação do Edital 003/2026 e preparação futura}
O Edital 003/2026 informa que o \textbf{Anexo III} contém os conteúdos programáticos e o \textbf{Anexo XVI} os pontos da prova didática. Na data deste levantamento, \textbf{não foi possível acessar e conferir integralmente o Anexo III de Matemática do concurso 2026}, tampouco identificar o caderno oficial da área em endereço verificável. Portanto, \emph{não existe aqui uma lista atribuída ao programa de 2026}. Verificar assim que disponíveis:
\begin{enumerate}
\item texto e eventuais retificações do Anexo III;
\item quantidade de questões e distribuição das disciplinas no caderno de Matemática;
\item caderno original, gabarito preliminar, recursos e gabarito definitivo;
\item temas efetivamente definidos no Anexo XVI.
\end{enumerate}
Página do concurso: \url{https://seletivo.ifmt.edu.br/edital/visualizar/173/}.

\section{Plano de estudo baseado nas duas provas verificadas}
\begin{enumerate}
\item Resolver integralmente os dois cadernos com tempo e sem consulta, mantendo registro de erros por assunto.
\item Para Matemática, organizar os 17 eixos programáticos; começar pelos tópicos com maior dificuldade e recorrência observada.
\item Estudar separadamente as partes comuns (Português, Fundamentos, MT) para evitar nota insuficiente por negligência de conteúdo geral.
\item Treinar elaboração de provas didáticas apenas depois de confirmar os temas sorteáveis do edital-alvo; reutilizar os planos já existentes no diretório de aulas como material de apoio.
\item Atualizar esta comparação quando o Anexo III de 2026 e o caderno específico forem localizados.
\end{enumerate}

\section*{Documentação oficial}
\begin{enumerate}
\item IFMT. Edital 166/2023, Anexo III consolidado. \url{https://seletivo.ifmt.edu.br/edital/arquivos/889/}.
\item IFMT. Edital 029/2024, Anexo III consolidado. \url{https://processoseletivo.ifmt.edu.br/media/filer_public/2c/de/2cde91b5-8ae1-4067-8e0c-37e67430f39f/edital_2024_029__anexos_cp_ebtt_consolidado_01.pdf}.
\item IFMT. Edital 003/2026, acompanhamento institucional. \url{https://seletivo.ifmt.edu.br/edital/visualizar/173/}.
\item Cadernos completos de Matemática, Editais 166/2023 e 029/2024, conforme links da seção 2.
\end{enumerate}
\end{document}
